\section{Cryptographic Voting and Tally Protocol}
\label{sec:05}
\label{sec:cryptographic-protocol}

\subsection{Notation and public statement}
\label{subsec:protocol-notation}

The protocol uses a custom Circom voting relation and an election-specific
\texttt{VoteChainDirect} contract. The circuit operates over the BN254
scalar field $\mathbb{F}$ and uses Poseidon for commitment and context calculations
\citep{grassi2021poseidon}. Ballot encryption uses the prime-order BabyJub
subgroup $\mathbb{G}$. Table~\ref{tab:protocol-notation} gives the notation used in
Sections~\ref{sec:system-model}--\ref{sec:cryptographic-protocol}.

% Keep the notation table with this section rather than allowing a float
% to move above the section heading in the CAS single-column class.
\begin{center}
\begin{minipage}{\linewidth}
\centering
\makeatletter\def\@captype{table}\makeatother
\caption{Consolidated notation for the private-voting and tally protocol.}
\label{tab:protocol-notation}
\footnotesize
\begin{tabular}{@{}p{0.29\linewidth}p{0.63\linewidth}@{}}
\toprule
\textbf{Symbol} & \textbf{Meaning} \\
\midrule
$e$, $s$ & Election identifier and private voter credential secret \\
$\mathbb{F}$, $\mathbb{G}$, $q$ & Circuit field, prime-order BabyJub subgroup, and subgroup order \\
$G$, $Y=(Y_x,Y_y)$ & Subgroup generator and election encryption public key \\
$d_{\mathrm{id}}$, $d_{\mathrm{null}}$, $d_{\mathrm{ctx}}$ & Poseidon domain separators for identity, nullifier, and election context \\
$H_P$ & Poseidon hash with the ordered inputs shown in each equation \\
$\rho_e$, $p$ & Frozen election eligibility root and private Merkle authentication path \\
$\mu_e$, $\delta$, $\chi_e$ & Candidate-manifest hash, contract-domain value, and election-context hash \\
$K$, $M$, $L$ & Active candidates, candidate-slot capacity ($16$), and Merkle depth ($17$) \\
$C_s$, $n_e$ & Voter identity commitment and election-scoped nullifier \\
$v_j$, $r_j$ & Private ballot bit and encryption scalar for slot $j$ \\
$C_j=(R_j,S_j)$ & Exponential-ElGamal ciphertext for ballot slot $j$ \\
$\mathbf{w}$, $\mathbf{u}$, $\pi$ & Private circuit witness, public-signal vector, and Groth16 proof \\
$\mathcal{R}_{\mathrm{VC}}$ & Constrained relation verified by the direct voting circuit \\
$B$, $\overline{C}_j$ & Number of accepted ballots and encrypted aggregate for slot $j$ \\
$x_i$, $X_i$ & Trustee $i$'s private share and corresponding public share $x_iG$ \\
$D_{i,j}$, $\gamma_{i,j}$ & Trustee $i$'s partial decryption and its consistency proof for slot $j$ \\
$I$, $\lambda_i$, $t_j$ & Valid trustee set, reconstruction coefficient, and recovered slot total \\
$T_{\mathrm{DKG}}$, $T_{\mathrm{fin}}$ & Distributed-key-generation and finalization transcripts \\
$h_{\mathrm{DKG}}$, $h_{\mathrm{fin}}$ & Canonical SHA-256 hashes of the corresponding transcripts \\
\bottomrule
\end{tabular}
\end{minipage}
\end{center}

The compiled circuit fixes $M=16$ candidate slots and a Merkle depth of
$L=17$. The active candidate count satisfies $1\leq K\leq M$.
These values describe the current implementation, rather than a variable-size
circuit family. The Merkle path has one sibling and one direction bit per
level. Its root is frozen in the election contract before voting begins.
The Groth16 verifier receives a public vector containing nine election and
nullifier values followed by four coordinate arrays of length $M$:
\begin{equation}
\label{eq:protocol-public-inputs}
\mathbf{u}=\bigl(
\rho_e,e,\mu_e,K,\delta,Y_x,Y_y,\chi_e,n_e,
\mathbf{R}_x,\mathbf{R}_y,\mathbf{S}_x,\mathbf{S}_y
\bigr).
\end{equation}
The resulting dimension is $9+4M=73$. The prover's private witness
$\mathbf{w}$ contains the credential secret, Merkle siblings and path
bits, candidate-selection bits, and encryption scalars. The circuit
does not take a civil identifier or biometric measurement as an input.

% Verified source:
% circuits/VoteChainVoteDirect.circom:11-143
% circuits/config/VoteChainBuildProfile.circom:13-54
% contracts/wp4-foundry/src/VoteChainWP4.sol:825-917

\subsection{Eligibility, nullifier, and election binding}
\label{subsec:protocol-eligibility}

The custom credential secret $s$ is a nonzero field element. Its identity
commitment is derived using the identity domain separator:
\begin{equation}
\label{eq:protocol-identity}
C_s=H_P(d_{\mathrm{id}},s).
\end{equation}
The circuit reconstructs the Merkle root from $C_s$ and the private
authentication path $p$. At every level, a Boolean direction bit determines
the order of the current node and its sibling. The final value must equal
the configured root $\rho_e$:
\begin{equation}
\label{eq:protocol-merkle}
\operatorname{MerkleRoot}_{L}(C_s,p)=\rho_e.
\end{equation}
The path therefore establishes membership of the committed credential in
the frozen eligibility set. It does not perform biometric recognition.
That earlier authorization step supplies access to the credential and path.

The circuit also enforces the election-scoped nullifier already defined in
Eq.~\eqref{eq:system-nullifier}. The same credential and election identifier
produce the same public nullifier. This permits the contract to reject a
second accepted ballot using that value. The mechanism does not prevent
multiple credentials from being issued to one person during enrollment.
The election-context hash binds the public election identifier, manifest,
active candidate count, contract-domain value, and encryption key:
\begin{equation}
\label{eq:protocol-context}
\chi_e=H_P\bigl(
 d_{\mathrm{ctx}},e,\mu_e,K,\delta,Y_x,Y_y
\bigr).
\end{equation}
The circuit recomputes this hash from the declared public values. The
contract supplies those values from its stored election configuration.
The voter-facing web client separately recomputes the candidate manifest
from the registered candidate identifiers and compares it with the
contract configuration. The circuit binds the manifest hash but does not
reconstruct the candidate list from individual candidate records.

% Verified source:
% circuits/lib/VoteChainVoteCore.circom:78-243
% votechain-frontend-ts/src/privacy/wp5Privacy.worker.ts:680-814
% contracts/wp4-foundry/src/VoteChainWP4.sol:825-864

\subsection{Ballot constraints and ciphertext construction}
\label{subsec:protocol-encryption}

The selected candidate is encoded as a private binary vector of length $M$.
The bit $v_j$ represents candidate slot $j$, with indices starting at zero.
The circuit constrains each bit, requires exactly one active selection, and
sets all padding slots to zero:
\begin{equation}
\label{eq:protocol-onehot}
\begin{aligned}
 v_j&\in\{0,1\} && (0\leq j<M),\\
 \sum_{j=0}^{M-1}v_j&=1,\\
 v_j&=0 && (K\leq j<M).
\end{aligned}
\end{equation}
It also checks the range $1\leq K\leq16$. These constraints exclude
empty, multiple-choice, and out-of-manifest plaintext ballot vectors
within the stated single-choice election model. Such proof checks for
exponential-ElGamal ballots are established prior art
\citep{huber2025ballotvalidity}.

For ballot encryption, $G$ is the configured BabyJub subgroup generator.
The public key $Y$ must lie on the curve, be non-identity, and belong to
the prime-order subgroup. The circuit checks these conditions. For each
slot, the client supplies a private scalar $r_j$ with $1\leq r_j<q$.
The corresponding exponential-ElGamal ciphertext is
\begin{equation}
\label{eq:protocol-elgamal}
\begin{aligned}
 R_j &= r_jG,\\
 S_j &= r_jY+v_jG,\\
 C_j &= (R_j,S_j),\qquad 0\leq j<M.
\end{aligned}
\end{equation}
Here, the group identity represents an encrypted zero contribution after
removing the shared-key term. Every slot is encrypted, including padding
slots. The client generates the encryption scalars locally. The circuit
checks their range and the ciphertext equations, not the entropy of the
client's random-number generator.
The direct circuit wrapper equates its 64 public ciphertext coordinates
with the points calculated by the core circuit. This equality is part of
the proof statement. Consequently, a valid proof cannot be paired with a
different ciphertext array without causing verification to fail, subject
to the proof system's soundness assumptions.

% Verified source:
% circuits/lib/VoteChainVoteCore.circom:245-488
% circuits/VoteChainVoteDirect.circom:89-143
% votechain-frontend-ts/src/privacy/wp5Privacy.worker.ts:855-1147

\subsection{Zero-knowledge proof and on-chain ballot acceptance}
\label{subsec:protocol-verification}

The direct voting relation combines credential membership, nullifier
consistency, election-context binding, one-hot validity, and correctness of
all encrypted slots. It can be expressed as
\begin{equation}
\label{eq:protocol-relation}
\mathcal{R}_{\mathrm{VC}}(\mathbf{u},\mathbf{w})=1
\end{equation}
when the constraints in Eqs.~\eqref{eq:protocol-identity}--
\eqref{eq:protocol-elgamal} hold for the public inputs in
Eq.~\eqref{eq:protocol-public-inputs}. The private path and credential
remain within the witness. The locally selected candidate is likewise
not included among the public signals.

The client runs the Groth16 prover with the voting circuit and its proving
key. It checks the returned public signals against the expected election
configuration and ciphertext coordinates. The client then verifies the
proof locally before preparing the transaction. This local check does not
replace on-chain verification \citep{groth2016snark}.
A call to \texttt{VoteChainDirect.castVote} supplies $\pi$, $n_e$, and the
64 ciphertext coordinates. The contract requires an active election and
an unused nullifier. It checks that the ciphertext coordinates are in the
canonical field range. It then reconstructs the 73 public inputs using its
stored configuration and calls the Groth16 verifier. Only after successful
verification does the contract record the nullifier and update the tally.

The proof attests to the circuit relation, not to a fresh biometric match.
Its validity also depends on the correct proving and verification artifacts.
The research implementation does not establish an independently audited,
production-grade trusted-setup ceremony.

% Verified source:
% votechain-frontend-ts/src/privacy/wp5Privacy.worker.ts:1027-1147
% contracts/wp4-foundry/src/VoteChainWP4.sol:825-976

\subsection{Encrypted aggregation and threshold decryption}
\label{subsec:protocol-threshold}

The contract initializes each encrypted candidate aggregate to the BabyJub
identity points. After accepting a ballot, it adds the corresponding
ciphertext points to every slot. For $B$ accepted ballots, the aggregate at
slot $j$ is
\begin{equation}
\label{eq:protocol-aggregation}
\overline{C}_j=(\overline{R}_j,\overline{S}_j)
 =\sum_{b=1}^{B}C_{b,j}.
\end{equation}
The addition is componentwise in the BabyJub group. By the ElGamal
construction, the resulting second point contains the encoded candidate
total $t_jG$ masked by the aggregate encryption term. The contract does
not maintain plaintext candidate counts during the active private ballot
phase.

Multi-authority encrypted election tallies have established precedent
\citep{cramer1997multiauthority}. The implementation specifies five
trustees and a reconstruction threshold of three. The verified distributed-key-generation
transcript defines the election public key and public trustee shares.
Each trustee holds a private share $x_i$ associated with its public point
$X_i=x_iG$. For an aggregate $\overline{C}_j$, a trustee can produce
\begin{equation}
\label{eq:protocol-partial}
D_{i,j}=x_i\overline{R}_j.
\end{equation}
The partial-decryption record contains a discrete-logarithm equality
proof $\gamma_{i,j}$. It binds the trustee identifier, election identifier,
slot number, public share, aggregate point, and partial decryption.
In the implemented non-interactive proof, the trustee forms commitment
points $A_{1,i,j}=k_{i,j}G$ and
$A_{2,i,j}=k_{i,j}\overline{R}_j$. A domain-separated Poseidon challenge
$c_{i,j}$ is computed from the election and trustee context together with
the relevant points. The response is
$z_{i,j}=k_{i,j}+c_{i,j}x_i\pmod q$.
The transcript verifier recomputes the challenge and checks
\begin{equation}
\label{eq:protocol-dleq}
\begin{aligned}
 z_{i,j}G&=A_{1,i,j}+c_{i,j}X_i,\\
 z_{i,j}\overline{R}_j&=A_{2,i,j}+c_{i,j}D_{i,j}.
\end{aligned}
\end{equation}
These checks establish consistency of an accepted partial decryption with
the public trustee share and the specified aggregate point, under the
proof's cryptographic assumptions. They are implemented in the external
transcript verifier rather than the voting contract.

For a set $I$ of at least three distinct trustees with valid records, the
verifier combines partial decryptions using Lagrange coefficients at zero:
\begin{equation}
\label{eq:protocol-lagrange}
\begin{aligned}
 \lambda_i&=\prod_{\substack{k\in I\\k\neq i}}
 \frac{-k}{i-k}\pmod q,\\
 D_j&=\sum_{i\in I}\lambda_iD_{i,j}.
\end{aligned}
\end{equation}
The recovered message point is
\begin{equation}
\label{eq:protocol-tally}
 t_jG=\overline{S}_j-D_j,\qquad 0\leq t_j\leq B.
\end{equation}
The verifier decodes $t_j$ within the accepted-ballot bound using a
bounded baby-step/giant-step search. It rejects a transcript if there are
fewer than three valid partial records for any slot. It also rejects
nonzero totals in padded slots and a final active-slot sum different from
$B$.

The threshold controls who can collectively decrypt. It does not make
individual encrypted ballots intrinsically undecryptable to a colluding
threshold of trustees. Confidentiality therefore depends on the actual
custody and independence of the private shares.

% Verified source:
% contracts/wp4-foundry/src/VoteChainWP4.sol:690-801
% crypto/src/wp6-threshold.mjs:27-43,1643-1719,1721-2040
% crypto/src/wp6-threshold.mjs:2042-2150,2307-2654

\subsection{Finalization transcript and audit boundary}
\label{subsec:protocol-finalization}

The finalization evidence contains the verified DKG transcript, encrypted
candidate aggregates, trustee partial-decryption records, and reconstructed
totals. The application first checks that the election identifiers, frozen
configuration, accepted-ballot count, and aggregate ciphertexts match the
corresponding on-chain values. The threshold verifier then validates the
transcript and recovers the candidate totals.

Both transcript commitments use canonical JSON serialization followed by
SHA-256. Their relation to the published evidence is
\begin{equation}
\label{eq:protocol-transcript-hash}
\begin{aligned}
 h_{\mathrm{DKG}}&=\mathsf{SHA256}
 (\mathsf{Canon}(T_{\mathrm{DKG}})),\\
 h_{\mathrm{fin}}&=\mathsf{SHA256}
 (\mathsf{Canon}(T_{\mathrm{fin}})).
\end{aligned}
\end{equation}
The DKG hash is pinned during election configuration. Before finalization,
the application requires the supplied DKG transcript to match this hash.
It also checks that the supplied aggregates match the contract state.
A verified finalization produces the plaintext totals and
$h_{\mathrm{fin}}$ for submission.

The administrator calls \texttt{finalizeTally} after the election enters its
Ended phase. The contract requires a previously pinned DKG-transcript hash
and a nonzero finalization hash. It checks that padded totals are zero,
that each active total is no larger than $B$, and that their sum equals
$B$. The contract then records the totals and $h_{\mathrm{fin}}$ and moves
to Finalized.

The on-chain transition does not verify the DKG transcript, DLEQ proofs,
or the derivation of plaintext totals from the encrypted aggregate. Those
cryptographic checks occur outside the contract. A complete audit therefore
requires the finalization transcript, the pinned hashes, and the published
encrypted aggregates. The recorded hash allows an auditor to detect a
mismatch between the supplied transcript and the on-chain commitment.
It does not force the administrator to submit only a cryptographically
verified transcript. The contract's arithmetic checks must not be
interpreted as an on-chain proof of correct threshold decryption.
The protocol thus separates two verifiable statements. Groth16 is checked
on-chain when an encrypted ballot is accepted. Trustee-decryption evidence
is checked by the external verifier during result disclosure. The next
section documents the corresponding implementation and experimental
procedures without treating these distinct checks as equivalent.

% Verified source, pinned revision:
% votechain-central-server/services/elections/electionFinalizationService.js:994-1257
% crypto/src/wp6-threshold.mjs:2307-2654
% contracts/wp4-foundry/src/VoteChainWP4.sol:532-616
% Git commit 17c7453c50793643b590fb5beac6ad77b9aa53ce